\documentclass[12pt,a4paper]{article}

% ---------- Кодировки ----------
\usepackage[T2A]{fontenc}
\usepackage[utf8]{inputenc}
\usepackage[russian,english]{babel}
\usepackage{cmap}

% ---------- Геометрия ----------
\usepackage[left=2.5cm,right=2cm,top=2cm,bottom=2cm]{geometry}
\usepackage{setspace}
\onehalfspacing

% ---------- Пакеты ----------
\usepackage{graphicx}
\usepackage{float}
\usepackage{hyperref}
\usepackage{enumitem}
\usepackage{booktabs}
\usepackage{amsmath,amssymb}
\usepackage{xcolor}

% ---------- Гиперссылки ----------
\hypersetup{colorlinks=true,linkcolor=blue!50!black,urlcolor=blue!50!black}

% ---------- Вставка скриншота ----------
\newcommand{\screenshot}[3]{%
    \IfFileExists{#1}{%
        \begin{figure}[H]
            \centering
            \includegraphics[width=0.9\textwidth]{#1}
            \caption{#2}
            \label{#3}
        \end{figure}
    }{}%
}

% ---------- Авторы ----------
\newcommand{\AuthorA}{Ясаков Артём}
\newcommand{\AuthorB}{Сарваров Тимур}
\newcommand{\Group}{P3313}
\newcommand{\Teacher}{Жданов Андрей Дмитриевич}

% ============================================================
\begin{document}

% ---------- Титульный лист ----------
\begin{titlepage}
    \centering
    {\large Федеральное государственное автономное образовательное\\
    учреждение высшего образования\\
    «Национальный исследовательский университет ИТМО»\par}
    {\large Факультет программной инженерии и компьютерной техники\par}
    \vspace{2.5cm}
    {\Huge\bfseries Лабораторная работа №6\par}
    \vspace{0.4cm}
    {\Large Логистическая регрессия\par}
    \vspace{0.4cm}
    {\Large Датасет Titanic\par}
    \vspace{4cm}
    \begin{flushright}
    \begin{tabular}{rl}
        \textbf{Выполнили:} & \AuthorA \\
                            & \AuthorB \\[6pt]
        \textbf{Группа:}    & \Group \\[12pt]
        \textbf{Проверил:}  & \Teacher \\
    \end{tabular}
    \end{flushright}
    \vfill
    {\large Санкт-Петербург, 2026 г.\par}
\end{titlepage}

\tableofcontents
\newpage

% ============================================================
\section{Цель работы}

Изучить метод логистической регрессии и реализовать его «с нуля»
на языке Python, используя только библиотеки \texttt{NumPy}
и \texttt{Pandas}. Провести предварительную обработку данных
датасета Titanic, визуализировать статистику, исследовать влияние
гиперпараметров (learning rate, количество итераций, метод
оптимизации) на качество модели. Оценить модель с помощью метрик
Accuracy, Precision, Recall и F1-Score.

% ============================================================
\section{Описание датасета}

Используется датасет \texttt{titanic} из библиотеки
\texttt{seaborn}. Он содержит информацию о пассажирах
трагически известного лайнера «Титаник», затонувшего в 1912 году.
Целевая переменная — выжил ли пассажир (\texttt{survived}:
0 — погиб, 1 — выжил).

\begin{itemize}[nosep]
    \item Количество образцов: \textbf{891}
    \item Количество исходных признаков: \textbf{15}
    \item Тип задачи: бинарная классификация
    \item Классы: погиб (0) — 549 чел. (61.6\%), выжил (1) — 342 чел. (38.4\%)
\end{itemize}

\subsection{Используемые признаки}

Из исходных столбцов были отобраны следующие признаки для моделирования:

\begin{itemize}[nosep]
    \item \texttt{pclass} — класс билета (1, 2, 3)
    \item \texttt{sex} — пол (\texttt{male} / \texttt{female})
    \item \texttt{age} — возраст (годы)
    \item \texttt{sibsp} — количество братьев/сестёр/супругов на борту
    \item \texttt{parch} — количество родителей/детей на борту
    \item \texttt{fare} — стоимость билета
    \item \texttt{embarked} — порт посадки (C = Cherbourg, Q = Queenstown, S = Southampton)
\end{itemize}

Столбцы \texttt{deck} (слишком много пропусков — 688), \texttt{alive}
(дублирует целевую), \texttt{class}, \texttt{who}, \texttt{adult\_male},
\texttt{alone} не использовались.

% ============================================================
\section{Предварительная обработка данных}

\subsection{Пропущенные значения}

В исходном датасете присутствуют пропуски:
\texttt{age} — 177 значений (19.9\%), \texttt{embarked} — 2 значения (0.2\%),
\texttt{deck} — 688 значений (77.2\%). Признак \texttt{deck} исключён из
анализа. Пропуски в \texttt{age} и \texttt{fare} заполнены медианным
значением соответствующего столбца: \(\tilde{x}_{\text{age}} = 28\),
\(\tilde{x}_{\text{fare}} = 14.45\). Пропуски в \texttt{embarked}
незначительны (2 шт.) и также заполнены модой.

\subsection{Кодирование категориальных признаков}

Признак \texttt{sex} преобразован в бинарный:
\(\text{female} = 1\), \(\text{male} = 0\).
Признаки \texttt{embarked} и \texttt{pclass} преобразованы с помощью
one-hot кодирования (\texttt{pd.get\_dummies}) с параметром
\texttt{drop\_first=True} для избежания мультиколлинеарности.
Итоговый набор — 9 числовых признаков:
\texttt{sex}, \texttt{age}, \texttt{sibsp}, \texttt{parch},
\texttt{fare}, \texttt{embarked\_Q}, \texttt{embarked\_S},
\texttt{pclass\_2}, \texttt{pclass\_3}.

\subsection{Нормализация}

Все признаки масштабированы с помощью \texttt{StandardScaler}:
\[
x' = \frac{x - \mu}{\sigma},
\]
где \(\mu\) — среднее, \(\sigma\) — стандартное отклонение признака.
Масштабирование необходимо для стабильной сходимости градиентного
спуска и корректной работы регуляризации.

\subsection{Разделение данных}

Датасет разделён на обучающую и тестовую выборки в соотношении
\textbf{80\,/\,20} со стратификацией по целевой переменной
(\texttt{stratify=y}). Итоговые размеры:
\[
\text{train} = 712 \text{ образцов},\quad \text{test} = 179 \text{ образцов}.
\]
Доля выживших: train — 38.3\%, test — 38.5\% (одинакова благодаря
стратификации).

% ============================================================
\section{Статистика и визуализация}

\subsection{Описательная статистика}

В таблице~\ref{tab:stats} приведена описательная статистика признаков
до масштабирования: количество, среднее, стандартное отклонение,
минимум, квартили, максимум.

\begin{table}[H]
    \centering
    \caption{Описательная статистика признаков (до масштабирования)}
    \label{tab:stats}
    \small
    \begin{tabular}{lrrrrrrr}
        \toprule
        Признак & count & mean & std & min & 25\% & 50\% & max \\
        \midrule
        sex       & 891 & 0.352 & 0.478 & 0 & 0 & 0 & 1 \\
        age       & 891 & 29.36 & 13.02 & 0.42 & 22.0 & 28.0 & 80.0 \\
        sibsp     & 891 & 0.523 & 1.103 & 0 & 0 & 0 & 8 \\
        parch     & 891 & 0.382 & 0.806 & 0 & 0 & 0 & 6 \\
        fare      & 891 & 32.20 & 49.69 & 0 & 7.91 & 14.45 & 512.33 \\
        embarked\_Q & 891 & 0.086 & 0.281 & 0 & 0 & 0 & 1 \\
        embarked\_S & 891 & 0.723 & 0.448 & 0 & 0 & 1 & 1 \\
        pclass\_2   & 891 & 0.207 & 0.405 & 0 & 0 & 0 & 1 \\
        pclass\_3   & 891 & 0.551 & 0.498 & 0 & 0 & 1 & 1 \\
        \bottomrule
    \end{tabular}
\end{table}

Из таблицы видно, что 72.3\% пассажиров сели в порту Southampton,
55.1\% путешествовали третьим классом, средний возраст — около 29 лет.
Признак \texttt{fare} имеет сильную правостороннюю асимметрию
(медиана 14.45, максимум — 512.33).

\subsection{Гистограммы признаков}

Распределения признаков представлены на рисунке~\ref{fig:hist}.
Вертикальными линиями отмечены среднее (красная) и медиана (зелёная).

\screenshot{feature_distributions.png}{Распределение признаков датасета Titanic}{fig:hist}

% ============================================================
\section{Реализация логистической регрессии}

Модель логистической регрессии реализована «с нуля» без использования
готовых реализаций \texttt{sklearn} (кроме \texttt{train\_test\_split}
и \texttt{StandardScaler}). Реализация включает два класса, два метода
оптимизации и самостоятельный расчёт метрик.

\subsection{Сигмоидная функция (гипотеза)}

Сигмоидная функция преобразует линейную комбинацию признаков
в вероятность принадлежности к классу 1:
\[
\sigma(z) = \frac{1}{1 + e^{-z}},\quad z = w^T x + b.
\]
Для численной устойчивости аргумент \(z\) обрезается до диапазона
\([-500, 500]\) перед вычислением экспоненты.

\subsection{Функция потерь (Log Loss / Binary Cross-Entropy)}

\[
\mathcal{L}(w, b) = -\frac{1}{n} \sum_{i=1}^{n}
\bigl[ y_i \ln \hat{y}_i + (1 - y_i) \ln (1 - \hat{y}_i) \bigr],
\]
где \(\hat{y}_i = \sigma(w^T x_i + b)\). Для предотвращения
\(\ln(0)\) вероятности обрезаются до \([10^{-15}, 1 - 10^{-15}]\).

\subsection{Градиентный спуск}

Градиенты функции потерь по параметрам:
\[
\frac{\partial \mathcal{L}}{\partial w} = \frac{1}{n} X^T (\hat{y} - y),
\qquad
\frac{\partial \mathcal{L}}{\partial b} = \frac{1}{n} \sum_{i=1}^{n} (\hat{y}_i - y_i).
\]
Правило обновления:
\[
w := w - \eta \frac{\partial \mathcal{L}}{\partial w},\quad
b := b - \eta \frac{\partial \mathcal{L}}{\partial b},
\]
где \(\eta\) — коэффициент обучения (learning rate). Класс —
\texttt{LogisticRegressionGD}.

\subsection{Метод Ньютона (Newton-Raphson)}

Метод Ньютона использует вторые производные (матрицу Гессе) для более
быстрой сходимости:
\[
H = \frac{1}{n} X_{\text{aug}}^T \, S \, X_{\text{aug}},
\]
где \(S = \operatorname{diag}\bigl(\hat{y}_i (1 - \hat{y}_i)\bigr)\),
\(X_{\text{aug}} = [X \mid \mathbf{1}]\).

Градиент и Гессиан объединяются для параметров \(w\) и \(b\):
\[
\theta := \theta - H^{-1} \nabla \mathcal{L}(\theta).
\]
Для устойчивости к матрице Гессе добавляется регуляризация:
\(H_{\text{reg}} = H + 10^{-5} I\). Класс — \texttt{LogisticRegressionNewton}.

\subsection{Метрики качества}

Все метрики реализованы вручную (функция \texttt{compute\_metrics}):

\[
\text{Accuracy} = \frac{TP + TN}{TP + TN + FP + FN},
\]
\[
\text{Precision} = \frac{TP}{TP + FP},
\quad
\text{Recall} = \frac{TP}{TP + FN},
\]
\[
F_1 = 2 \cdot \frac{\text{Precision} \cdot \text{Recall}}{\text{Precision} + \text{Recall}}.
\]

% ============================================================
\section{Исследование гиперпараметров}

Проведено исследование влияния трёх гиперпараметров на качество модели:
коэффициента обучения \(\eta\), количества итераций \(N\) и метода
оптимизации (градиентный спуск vs метод Ньютона).

\subsection{Влияние learning rate}

Исследованы значения \(\eta \in \{0.001, 0.01, 0.05, 0.1, 0.5, 1.0\}\)
при фиксированном \(N = 5000\). Результаты приведены в
таблице~\ref{tab:lr}.

\begin{table}[H]
    \centering
    \caption{Влияние learning rate (GD, \(N = 5000\))}
    \label{tab:lr}
    \begin{tabular}{cccccc}
        \toprule
        \(\eta\) & Accuracy & Precision & Recall & F1-Score \\
        \midrule
        0.001 & 0.7765 & 0.7302 & 0.6667 & 0.6970 \\
        0.01  & 0.7989 & 0.7797 & 0.6667 & 0.7188 \\
        0.05  & 0.8045 & 0.7931 & 0.6667 & 0.7244 \\
        0.1   & \textbf{0.8045} & \textbf{0.7931} & 0.6667 & \textbf{0.7244} \\
        0.5   & 0.8045 & 0.7931 & 0.6667 & 0.7244 \\
        1.0   & 0.8045 & 0.7931 & 0.6667 & 0.7244 \\
        \bottomrule
    \end{tabular}
\end{table}

При \(\eta = 0.001\) модель сходится медленно. При \(\eta \geq 0.05\)
результаты стабилизируются. Расходимость при больших \(\eta\) на данном
датасете не наблюдается благодаря нормализации.

\subsection{Влияние количества итераций}

Исследованы \(N \in \{100, 500, 1000, 5000, 10000\}\) при \(\eta = 0.1\).
Результаты в таблице~\ref{tab:niter}.

\begin{table}[H]
    \centering
    \caption{Влияние количества итераций (GD, \(\eta = 0.1\))}
    \label{tab:niter}
    \begin{tabular}{cccccc}
        \toprule
        \(N\) & Accuracy & Precision & Recall & F1-Score \\
        \midrule
        100   & 0.7989 & 0.7797 & 0.6667 & 0.7188 \\
        500   & 0.7989 & 0.7797 & 0.6667 & 0.7188 \\
        1000  & 0.8045 & 0.7931 & 0.6667 & 0.7244 \\
        5000  & 0.8045 & 0.7931 & 0.6667 & 0.7244 \\
        10000 & \textbf{0.8045} & \textbf{0.7931} & 0.6667 & \textbf{0.7244} \\
        \bottomrule
    \end{tabular}
\end{table}

Модель выходит на плато уже при \(N = 1000\), дальнейшее увеличение
итераций не даёт существенного прироста качества.

\subsection{Метод Ньютона}

Исследованы \(N \in \{3, 5, 10, 15, 20, 30, 50\}\). Результаты
приведены в таблице~\ref{tab:newton}.

\begin{table}[H]
    \centering
    \caption{Результаты метода Ньютона}
    \label{tab:newton}
    \begin{tabular}{cccccc}
        \toprule
        \(N\) & Accuracy & Precision & Recall & F1-Score \\
        \midrule
        3  & 0.8045 & 0.7931 & 0.6667 & 0.7244 \\
        5  & 0.8045 & 0.7931 & 0.6667 & 0.7244 \\
        10 & 0.8045 & 0.7931 & 0.6667 & 0.7244 \\
        15 & 0.8045 & 0.7931 & 0.6667 & 0.7244 \\
        20 & 0.8045 & 0.7931 & 0.6667 & 0.7244 \\
        30 & 0.8045 & 0.7931 & 0.6667 & 0.7244 \\
        50 & 0.8045 & 0.7931 & 0.6667 & 0.7244 \\
        \bottomrule
    \end{tabular}
\end{table}

Метод Ньютона сходится к оптимальным параметрам уже за \textbf{3 итерации},
тогда как градиентному спуску требуется \(\sim 1000\) итераций. При этом
финальные метрики идентичны — точность не зависит от метода оптимизации
(оба сходятся к одному и тому же оптимуму).

\subsection{Визуализация}

На рисунке~\ref{fig:hyper} представлены графики зависимости метрик от
гиперпараметров, а на рисунке~\ref{fig:heatmap} — тепловая карта
Accuracy для всех комбинаций \(\eta\) и \(N\) (градиентный спуск).

\screenshot{hyperparameter_study.png}{Зависимость метрик от гиперпараметров}{fig:hyper}
\screenshot{heatmap_summary.png}{Тепловая карта Accuracy и сводная таблица}{fig:heatmap}

% ============================================================
\section{Финальная модель}

С оптимальными гиперпараметрами (\(\eta = 0.1\), \(N = 10000\),
градиентный спуск) получены следующие результаты на тестовой выборке:

\begin{table}[H]
    \centering
    \caption{Финальные метрики модели}
    \label{tab:final}
    \begin{tabular}{lcc}
        \toprule
        Метрика & Значение \\
        \midrule
        Accuracy  & 0.8045 \\
        Precision & 0.7931 \\
        Recall    & 0.6667 \\
        F1-Score  & 0.7244 \\
        \bottomrule
    \end{tabular}
\end{table}

Матрица ошибок:
\[
\begin{bmatrix}
\text{TN} = 98 & \text{FP} = 12 \\
\text{FN} = 23 & \text{TP} = 46
\end{bmatrix}
\]

Модель правильно классифицирует 98 из 110 погибших (Recall для класса
«погиб» \(\approx 0.89\)), но лишь 46 из 69 выживших (Recall для класса
«выжил» \(\approx 0.67\)). Это означает, что модель чаще ошибается
в~сторону предсказания «погиб», что консервативно и соответствует
априорному распределению классов (61.6\% — погибшие).

\screenshot{final_model.png}{Кривая обучения и матрица ошибок}{fig:final}

% ============================================================
\section{Анализ результатов}

\begin{itemize}
    \item \textbf{Learning rate:} при \(\eta = 0.001\) модель сходится
          слишком медленно (Accuracy \(\approx 0.78\)). Оптимальный
          диапазон: \(\eta \in [0.01, 0.1]\). При \(\eta \geq 0.5\)
          расходимости не происходит благодаря нормализации, но на
          более сложных датасетах возможна нестабильность.
    \item \textbf{Количество итераций:} при \(N = 100\) модель уже даёт
          Accuracy \(\approx 0.80\). После \(N = 1000\) наступает плато,
          дальнейшее обучение не улучшает метрики. Рекомендация:
          \(N = 5000\text{--}10000\) для надёжности.
    \item \textbf{Метод оптимизации:} метод Ньютона сходится за 3–5
          итераций против \(\sim 1000\) для градиентного спуска, но
          требует вычисления и обращения матрицы Гессе размера
          \((d+1)\times(d+1)\) (\(O(d^3)\)). Для данного датасета
          (\(d = 9\)) разница в скорости несущественна.
    \item \textbf{Precision vs Recall:} Precision (0.79) выше, чем Recall
          (0.67). Модель осторожна с предсказанием класса «выжил» —
          предпочитает пропустить выжившего, чем ошибочно предсказать
          выживание погибшему.
    \item \textbf{Общее качество:} Accuracy = 0.8045 — хороший результат
          для базовой логистической регрессии без feature engineering.
          Результат сопоставим с известными baseline для Titanic.
          Качество можно улучшить добавлением взаимодействий признаков
          (например, Age \(\times\) Pclass), полиномиальных признаков
          и~регуляризации L2.
\end{itemize}

% ============================================================
\section{Выводы}

В ходе лабораторной работы:

\begin{enumerate}
    \item Загружен датасет Titanic (891 пассажир, 15 исходных признаков),
          проведена предварительная обработка: отбор признаков (7
          содержательных), заполнение пропусков медианой, one-hot
          кодирование категориальных переменных, нормализация
          StandardScaler. Итоговый набор — 9 числовых признаков.
    \item Рассчитана и визуализирована описательная статистика:
          средние, стандартные отклонения, квартили, гистограммы
          распределений с наложением среднего и медианы.
    \item Данные разделены на обучающую (80\%) и тестовую (20\%) выборки
          со стратификацией по целевой переменной.
    \item Реализована логистическая регрессия «с нуля» на \texttt{NumPy}
          и \texttt{Pandas}: сигмоида, log loss, градиентный спуск,
          метод Ньютона. Метрики Accuracy, Precision, Recall, F1-Score
          реализованы вручную.
    \item Проведено исследование гиперпараметров: learning rate
          \(\{0.001, \dots, 1.0\}\), количество итераций
          \(\{100, \dots, 10000\}\), метод оптимизации (GD и Newton).
    \item Наилучшие результаты: Accuracy = \textbf{0.8045},
          F1-Score = \textbf{0.7244} при \(\eta = 0.1\), \(N = 10000\).
          Метод Ньютона достигает тех же метрик за 3 итерации.
\end{enumerate}

\newpage
\section*{Использование искусственного интеллекта}
\addcontentsline{toc}{section}{Использование искусственного интеллекта}

В ходе выполнения лабораторной работы применялись инструменты
искусственного интеллекта в качестве вспомогательного
средства. Конкретные сценарии использования:

\begin{enumerate}
    \item \textbf{Разбор задания.} Уточнение требований к реализации
          логистической регрессии «с нуля», интерпретация пунктов
          про исследование гиперпараметров и два метода оптимизации
          (градиентный спуск и метод Ньютона).
    \item \textbf{Реализация метода Ньютона.} Помощь в корректной
          реализации вычисления матрицы Гессе, её регуляризации и
          решения системы линейных уравнений для обновления параметров.
    \item \textbf{Численная устойчивость.} Рекомендации по обрезанию
          аргумента сигмоиды до \([-500, 500]\) и вероятностей до
          \([10^{-15}, 1 - 10^{-15}]\) для предотвращения переполнения.
    \item \textbf{Пояснение этапов обработки данных.} Объяснение
          необходимости one-hot кодирования с \texttt{drop\_first},
          нормализации для сходимости градиентного спуска и роли
          стратификации при разбиении.
    \item \textbf{Оформление отчёта.} Помощь в структурировании
          разделов, оформлении таблиц с результатами экспериментов
          и вставке графиков средствами \LaTeX{}.
    \item \textbf{Анализ результатов.} Интерпретация метрик Precision
          и Recall в контексте несбалансированности классов,
          обсуждение преимуществ и недостатков метода Ньютона по
          сравнению с градиентным спуском.
\end{enumerate}

% ============================================================
\end{document}